\documentclass[11pt]{article}
\usepackage[T1]{fontenc}
\usepackage{amsmath,amssymb,amsthm}
\usepackage[margin=1in]{geometry}
\usepackage[hidelinks,pdfusetitle]{hyperref}
\usepackage{microtype}

\newtheorem{theorem}{Theorem}
\newtheorem{lemma}[theorem]{Lemma}

\newcommand{\Nzero}{\mathbb N_0}
\newcommand{\TP}{T_{\mathrm P}}
\newcommand{\TL}{T_{\mathrm L}}

\title{No New Construction Is Needed for the Product and
\L{}ukasiewicz Cases of a Fuzzy-Metric Problem}
\author{Xiaojian Zeng\\
PBC School of Finance, Tsinghua University, 100190\\
\texttt{xiaojianzeng@foxmail.com}}
\date{}

\begin{document}
\maketitle

\begin{abstract}
Xu, Ji, and Qiu constructed a George--Veeramani fuzzy metric on
$\Nzero$ for the minimum t-norm and exhibited the sequence $x_n=n$ as
strictly fuzzy $\psi$-contractive but not Cauchy.  They asked whether such a
sequence exists for the product or the \L{}ukasiewicz t-norm.  We prove that
their function, sequence, and control function work unchanged for both
t-norms.  The decisive point is the direction of the fuzzy triangle
inequality: each requested t-norm is bounded above by the minimum t-norm.
For completeness, we verify both operations as continuous t-norms, all five
fuzzy-metric axioms, the exact contractivity identity, and a quantified
non-Cauchy witness.
\end{abstract}

\section{Context and statement}

George and Veeramani introduced the fuzzy-metric framework used here
\cite{GeorgeVeeramani1994}; standard background on continuous triangular
norms may be found in \cite{KlementMesiarPap2000}.  Gregori, Mi\~nana, and
Miravet introduced strictly fuzzy contractive sequences and asked whether
they must be Cauchy \cite{GregoriMinanaMiravet2020}.  Xu, Ji, and Qiu answered
that question negatively for a particular fuzzy metric equipped with the
minimum t-norm, and then asked for examples with the product and
\L{}ukasiewicz t-norms \cite[Problem~3]{XuJiQiu2026}.

The contribution below is source-local: no alternative fuzzy metric and no
alternative sequence are introduced.  Its proof is also an exact certificate,
with no numerical separation or search.  For comparison, explicit-certificate
arguments in other inequality problems can require several exact rational
separations to validate a proposed counterexample
\cite{ZengLiuRatnavelu2026}; here the entire transfer is instead forced by an
order relation between t-norms.

Let $\Nzero=\{0,1,2,\ldots\}$, and for $x,y\in\Nzero$ write
\[
 \ell(x,y)=\min\{x,y\},\qquad h(x,y)=\max\{x,y\}.
\]
For $t>0$, define
\[
 \mu_t(x,y)=\frac{|x-y|}{2^t-1}-\ell(x,y),
 \qquad
 q_t(x,y)=
 \begin{cases}
  0,&x=y,\\
  2^{\mu_t(x,y)-t},&x\ne y,
 \end{cases}
\]
and
\[
 M(x,y,t)=\frac{1}{1+q_t(x,y)}.
\]
The two operations under consideration are
\[
 \TP(a,b)=ab,
 \qquad
 \TL(a,b)=\max\{0,a+b-1\}
 \quad (a,b\in[0,1]).
\]

Recall that $(X,M,T)$ is a George--Veeramani fuzzy metric space if $T$ is a
continuous t-norm and, for all $x,y,z\in X$ and $s,t>0$,
\begin{align*}
 &M(x,y,t)>0, \tag{GV1}\\
 &M(x,y,t)=1\quad\Longleftrightarrow\quad x=y, \tag{GV2}\\
 &M(x,y,t)=M(y,x,t), \tag{GV3}\\
 &T\bigl(M(x,y,t),M(y,z,s)\bigr)\le M(x,z,t+s), \tag{GV4}\\
 &M(x,y,\mathord\cdot):(0,\infty)\longrightarrow(0,1]
     \text{ is continuous}. \tag{GV5}
\end{align*}
The orientation in (GV4) will be essential.

\begin{theorem}\label{thm:main}
For each $T\in\{\TP,\TL\}$, the triple $(\Nzero,M,T)$ is a
George--Veeramani fuzzy metric space.  Moreover, for
\[
 x_n=n\quad(n\in\mathbb N),
 \qquad
 \psi(u)=\frac{2u}{1+u}\quad(0<u\le1),
\]
the sequence $(x_n)$ is strictly fuzzy $\psi$-contractive and is not Cauchy.
Thus the same objects used by Xu, Ji, and Qiu for the minimum t-norm answer
their Problem~3 affirmatively for both requested t-norms.
\end{theorem}

\section{The triangle inequality}

We first rederive the stronger minimum-t-norm inequality, including the
degenerate cases needed below.

\begin{lemma}\label{lem:mintriangle}
For all $x,y,z\in\Nzero$ and $s,t>0$,
\[
 \min\{M(x,y,t),M(y,z,s)\}\le M(x,z,t+s).
\]
\end{lemma}

\begin{proof}
For $u,v\in\Nzero$, $r>0$, and $c\in\mathbb R$, multiplication by the
positive number $2^r-1$ gives the equivalence
\begin{equation}\label{eq:equivalence}
 \mu_r(u,v)\le c
 \quad\Longleftrightarrow\quad
 h(u,v)+c\le 2^r\bigl(\ell(u,v)+c\bigr).
\end{equation}
Fix $x,y,z,s,t$ and set
\[
 c=\max\{\mu_t(x,y),\mu_s(y,z)\}.
\]
Since $c\ge\mu_t(x,y)$, \eqref{eq:equivalence} gives
\[
 h(x,y)+c\le2^t\bigl(\ell(x,y)+c\bigr).
\]
Furthermore, $\ell(x,y)+c\ge0$: if $x\ne y$, it is at least
$|x-y|/(2^t-1)>0$, while if $x=y$, it is at least
$x+\mu_t(x,x)=0$.  Because $h(x,y)+c\ge\ell(x,y)+c\ge0$ and $2^t>1$,
we also have
\[
 \ell(x,y)+c\le2^t\bigl(h(x,y)+c\bigr).
\]
Irrespective of the order of $x$ and $y$, these two inequalities say
\[
 x+c\le2^t(y+c),\qquad y+c\le2^t(x+c).
\]
The same argument for $(y,z,s)$ gives
\[
 y+c\le2^s(z+c),\qquad z+c\le2^s(y+c).
\]
Composing in both directions and using $2^t2^s=2^{t+s}$ yields
\[
 h(x,z)+c\le2^{t+s}\bigl(\ell(x,z)+c\bigr).
\]
Another application of \eqref{eq:equivalence} therefore gives
\begin{equation}\label{eq:mu-triangle}
 \mu_{t+s}(x,z)\le
 \max\{\mu_t(x,y),\mu_s(y,z)\}.
\end{equation}

If $x,y,z$ are pairwise distinct, subtracting $t+s$ in
\eqref{eq:mu-triangle}, using $s,t>0$, and exponentiating gives
\begin{equation}\label{eq:qtriangle}
 q_{t+s}(x,z)\le\max\{q_t(x,y),q_s(y,z)\}.
\end{equation}
The remaining cases satisfy the same inequality.  If $x=z$, its left side
is zero.  If $x=y\ne z$, it reduces to
$q_{t+s}(x,z)\le q_s(x,z)$; this holds because, for a fixed off-diagonal
pair, the exponent
\[
 r\longmapsto \frac{|x-z|}{2^r-1}-\ell(x,z)-r
\]
is strictly decreasing on $(0,\infty)$.  The case $y=z\ne x$ is symmetric,
and the all-equal case is immediate.  Finally,
$q\mapsto(1+q)^{-1}$ is decreasing, so \eqref{eq:qtriangle} is precisely the
claimed inequality.
\end{proof}

\section{Proof of the theorem}

\begin{proof}[Proof of Theorem~\ref{thm:main}]
We first check the operations rather than merely naming them.  Multiplication
maps $[0,1]^2$ into $[0,1]$, is commutative, associative, coordinatewise
nondecreasing, continuous, and has neutral element $1$.  Thus $\TP$ is a
continuous t-norm.  The operation $\TL$ also maps $[0,1]^2$ into $[0,1]$,
is commutative, coordinatewise nondecreasing, continuous, and satisfies
$\TL(a,1)=a$.  For associativity, both bracketings of $a,b,c\in[0,1]$
equal
\[
 \max\{0,a+b+c-2\};
\]
indeed, if the inner sum in a bracketing is at most $1$, both that
bracketing and the displayed expression are zero, and otherwise direct
substitution gives the display.  Hence $\TL$ is a continuous t-norm as well.
Their endpoint behavior is $\TP(0,b)=\TL(0,b)=0$ and
$\TP(1,b)=\TL(1,b)=b$.

For every $a,b\in[0,1]$, including the endpoints,
\begin{equation}\label{eq:tnormorder}
 ab\le\min\{a,b\},
 \qquad
 \max\{0,a+b-1\}\le\min\{a,b\}.
\end{equation}
The first pair of inequalities follows from $a,b\le1$.  For the second,
$a+b-1\le a$ follows from $b\le1$, and $a+b-1\le b$ follows from
$a\le1$; zero is also at most both inputs.

We next verify all five fuzzy-metric axioms.  Since $t>0$, the denominator
$2^t-1$ is positive.  By definition $q_t(x,y)\ge0$, with equality exactly
on the diagonal.  Therefore $M(x,y,t)$ belongs to $(0,1]$ and equals $1$
exactly when $x=y$, proving (GV1) and (GV2).  Symmetry of $|x-y|$ and
$\ell(x,y)$ proves (GV3).  For fixed $x,y$, the function of $t$ is constant
equal to $1$ on the diagonal; off the diagonal it is a composition of
continuous functions on $(0,\infty)$.  This proves (GV5).

For (GV4), the two inputs to either t-norm belong to $(0,1]\subset[0,1]$.
Combining \eqref{eq:tnormorder} with Lemma~\ref{lem:mintriangle}, for either
$T\in\{\TP,\TL\}$ we obtain
\[
 T\bigl(M(x,y,t),M(y,z,s)\bigr)
 \le \min\{M(x,y,t),M(y,z,s)\}
 \le M(x,z,t+s).
\]
This is the required direction of (GV4), so all five axioms hold separately
under both labels.

It remains to check the sequence notions from their definitions.  Let
$\Psi$ be the set of continuous nondecreasing maps
$\psi:(0,1]\to(0,1]$ satisfying $\psi(u)>u$ for $0<u<1$.  The displayed
$\psi(u)=2u/(1+u)$ maps $(0,1]$ into itself, is continuous and increasing,
and
\[
 \psi(u)-u=\frac{u(1-u)}{1+u}>0\qquad(0<u<1).
\]
Thus $\psi\in\Psi$.

Shifting both arguments by one preserves their distance and increases their
minimum by one.  Off the diagonal this subtracts one from $\mu_t$, while on
the diagonal both $q$-values are zero.  Hence, in every case,
\begin{equation}\label{eq:shift}
 q_t(x+1,y+1)=\frac12q_t(x,y).
\end{equation}
Using $x_n=n$, equation \eqref{eq:shift} yields, for every
$m,n\in\mathbb N$ and $t>0$,
\begin{align*}
 M(x_{m+1},x_{n+1},t)
 &=\frac{1}{1+q_t(x_m,x_n)/2}\\
 &=\frac{2}{2+q_t(x_m,x_n)}\\
 &=\psi\bigl(M(x_m,x_n,t)\bigr).
\end{align*}
This equality implies the defining inequality for strict fuzzy
$\psi$-contractivity.  Neither the definition nor the identity evaluates
the selected t-norm.

Finally, a sequence is Cauchy here if for every $\varepsilon\in(0,1)$ and
$t>0$ there is $N\in\mathbb N$ such that
$M(x_n,x_m,t)>1-\varepsilon$ whenever $m,n\ge N$.  Given any proposed $N$,
take $n=N$, $m=2N+1$, $t=1$, and $\varepsilon=1/2$.  Then
\[
 \ell(x_N,x_{2N+1})=N,\qquad |x_N-x_{2N+1}|=N+1,
\]
so $\mu_1(x_N,x_{2N+1})=1$, $q_1(x_N,x_{2N+1})=1$, and
\[
 M(x_N,x_{2N+1},1)=\frac12.
\]
Both indices are at least $N$, but the Cauchy definition requires the strict
inequality $M>1/2$.  Thus the sequence is not Cauchy.  This definition also
contains no t-norm evaluation, so the same witness applies to both spaces.
\end{proof}

\section{Scope and assistance disclosure}

The result resolves the literal existence question in Problem~3 of the cited
version of Xu--Ji--Qiu by reusing its exact function, sequence, and control
function.  It does not assert a classification of all continuous t-norms or
of all strictly fuzzy contractive sequences, and it makes no priority claim
beyond this comparison with the cited sources.

An OpenAI Codex model assisted in organizing the proof, checking algebraic
case divisions, and preparing the manuscript.  A separate model pass
recomputed the decisive inequalities as an internal mathematical audit.  A
finite floating-point stress test of triangle cases was used only as a
diagnostic and is not evidence for any quantified statement.  Every theorem
in this note follows from the exact displayed argument.  The author reviewed
the final text and takes responsibility for its claims.

\bibliographystyle{plain}
\bibliography{references}

\end{document}
